% =====================================================================
% Introduction.
% Adapted from RESEARCH_PLAN.md §1. Tighter, third-person, present tense.
% =====================================================================
\section{Introduction}
\label{sec:intro}

A retrieval-augmented generation (RAG) prompt is a concatenation of retrieved
documents followed by the dialog:
\begin{equation*}
\underbrace{[\mathrm{doc}_1][\mathrm{doc}_2] \cdots [\mathrm{doc}_N]}_{\text{retrieved evidence}}
\; [\mathrm{query}]\;[\mathrm{answer}].
\end{equation*}
Standard causal attention makes prefill quadratic in the length of the
entire concatenation. Retrieved documents are often selected independently,
but whether cross-document interactions matter is an empirical question. The
prefill step still pays $O(T^2)$ for all such interactions.

Two costs follow, and the second dominates in production:
\begin{enumerate}
    \item \textbf{Compute.} Cross-document attention becomes the dominant term
          once $N$ grows.
    \item \textbf{Cache invalidation.} Under lower-triangular attention,
          $\mathrm{doc}_2$'s keys and values depend on $\mathrm{doc}_1$, so
          no document's KV block is reusable. Changing a single document in
          the retrieved set --- or even its position --- invalidates the
          entire prefix. A document that appears in a thousand requests is
          prefilled a thousand times.
\end{enumerate}
Prefix caching~\citep{kwon2023pagedattention} mitigates the second cost only
when \emph{exact} prefixes repeat; it does not admit the more common case in
which the same document appears at a different position or in a different
retrieved set.

\paragraph{Contribution.}
We characterize a mask family that isolates each retrieved document from
every other. The two design points explored in the prior literature ---
Fusion-in-Decoder~\citep{izacard2021fid} and Parallel Context
Windows~\citep{ratner2023pcw} --- correspond to a single choice on this axis;
this preprint measures three axes of that choice on single-hop
retrieval. Concretely:
\begin{itemize}
    \item We formalize the mask family (Section~\ref{sec:method}) and its
          cost model.
    \item We report descriptive full-baseline seed intervals and use their
          spread as context rather than as mask-comparison inference
          (Section~\ref{sec:method}).
    \item We measure quality changes after removing cross-document attention
          on the reported retrieval-QA cells
          (Section~\ref{sec:results-quality}) and report the realized
          prefill speedup on the same axes
          (Section~\ref{sec:results-efficiency}).
    \item We characterize cache correctness and a sequential-layout cache
          throughput estimate, and report a bounded restart-layout generation
          control (Section~\ref{sec:results-caching}).
\end{itemize}
The framing throughout is a characterization study rather than a benchmark
win. Self-distillation finetuning is evaluated for the train-free quality
gap on Qwen2.5-7B and Llama-3.1-8B using fixed held-out manifests and three
independently trained seeds per architecture. The sole admissible natural-QA
cross-architecture contrast does not establish a practical architecture
difference; unavailable cells are not interpreted as null effects. Direct comparison against PCW and
broader failure-envelope evaluation are out of scope for this preprint
(Section~\ref{sec:limitations}).
